\subsection{Channel Strip-Mode}\label{channelstrip}

The Channel Strip-Mode can be selected by pressing the \track button of the
VPOT ASSIGN section. It is built on the same base as the
\hyperref[panmode]{Pan-Mode}, so the \faders, the \select, \mute, \solo
and \rec buttons, the \bankdu/\channeldu navigation and the LEDs all work
exactly as described for the Pan-Mode. The only difference is that the
\vpots no longer act on the track assigned to the individual channel strip,
but on the \textbf{selected track}: all \vpots of a unit drive the
parameters of one plugin of the selected track's FX chain.

The idea behind the mode is to have quick access to the most used parameters
of a plugin without opening the FX window. To structure this, the mode works
on two levels:

\begin{itemize}
\item \textbf{Channel strips.} There are \textbf{16 global channel
  strips} that exist independently of any track or project. Each strip is
  bound to \emph{one} plugin and stores the mapping of that plugin's
  parameters onto 16 \vpots (VPOTs 1\,--\,8, and 9\,--\,16 with the \shift
  modifier). The strips are shared across all tracks and all projects and
  are stored in a single XML file in the user's Channel Strip Maps directory
  (see \hyperref[stripsave]{Saving and Loading Strips}).
\item \textbf{Assignments.} For every track and every hardware unit you can
  assign one of the 16 global strips (or none). The 8 \vpots of that unit
  (16 with \shift) then drive the 16 parameters of the assigned strip on the
  selected track. Two units of the same track can be assigned to different
  strips, e.g.~one for the EQ and one for the compressor. These assignments
  are stored in the \reaper project file.
\end{itemize}

The mode requires exactly \textbf{one selected track}. If no track, or more
than one track, is selected, no strip is active: the \vpots are dark and the
display shows the hint ``You must select a single track.''

\begin{itemize}
\item \faders: Control the track volume (or the pan, when \flip is active),
  as in the Pan-Mode.
\item \vpots: Control the parameters of the strip assigned to this unit.
  Turning a \vpot nudges the parameter; pressing a \vpot toggles the
  parameter between its two extreme values (e.g.~on/off).
\item \shift \vpots: Address the second 8 parameters of the strip (VPOTs
  9\,--\,16).
\item \select, \mute, \solo, \rec: As in the
  \hyperref[panmode]{Pan-Mode}.
\item \bankdu / \channeldu: As in the \hyperref[panmode]{Pan-Mode}.
\item \flip: Flip the function of the \faders and \vpots, as in the
  Pan-Mode.
\item \mfader: Control the master volume.
\end{itemize}

\subsubsection{Selecting and re-picking a strip}\label{stripselect}

How a unit's \vpots behave depends on whether a strip is assigned to that
unit for the selected track, and whether the plugin of that strip is present
on the track:

\begin{itemize}
\item \textbf{No strip (or plugin missing).} The display of the unit shows
  the names of the 16 global strips: VPOT 1\,--\,8 address strips
  1\,--\,8, and with \shift the VPOTs address strips 9\,--\,16. A strip whose
  plugin is not yet present on the selected track is marked with a leading
  ``\verb|+|''. Pressing a \vpot assigns the corresponding strip to this unit
  for the selected track and, if the plugin is still missing, adds it to the
  FX chain automatically (at the strip's insert position, see
  \hyperref[channelstripeditor]{next section}). If the same plugin is already
  on the track, the existing instance is reused and no second instance is
  added.
\item \textbf{Strip with plugin present.} The \vpots control the parameters
  of the assigned strip. The display shows the parameter name while idle and,
  shortly after turning a \vpot, its current value.
\end{itemize}

To change the strip of a unit that already has one, \textbf{press and hold
the \track button}. While it is held, every unit shows the strip-name list
again and a \vpot press re-assigns the strip for that unit. When you release
the \track button, the \vpots return to parameter control.

\subsubsection{Editor}\label{channelstripeditor}
The channel strips are edited in the \hyperref[F:Screenshot_ChannelStrip_Mode]{Channel
Strip-Mode editor}, which can be opened with \alt \track (see
Figure~\ref{F:Screenshot_ChannelStrip_Mode}). The editor is global: it shows
the 16 channel strips independently of the selected track or unit.

\pikFigure[0.62][11.5cm][-2.4mm]{Screenshot_ChannelStrip_Mode}{Channel
  Strip-Mode editor (\alt \track)}{}{}

Each row of the table represents one of the 16 global strips. The columns
are:

\begin{itemize}
\item \textbf{VPOT}: The position of the strip in the strip list (1\,--\,8
  and Shift 1\,--\,Shift 8). This is the number that the \vpots of a unit
  refer to when they show the strip-name list.
\item \textbf{Plugin}: The plugin bound to this strip. You can type to filter
  the list of installed plugins or pick one from the list. Selecting
  ``--- (no plugin) ---'' clears the whole strip (plugin, abbreviation and
  the \vpot mapping). When you pick a different plugin, the abbreviation is
  regenerated from the new plugin name.
\item \textbf{Abbrev}: A short display abbreviation (up to 5 characters)
  that is shown on the MCU. It is filled automatically from the plugin name
  and can be edited.
\item \textbf{Insert Position}: Where the plugin is inserted into the FX
  chain when a \vpot press adds it: at the \emph{first} position, at a fixed
  position 2\,--\,8, or \emph{last}.
  
  \par With REAPER 7.75 and the corresponding mixer option for empty FX
  slots, a fixed position (2\,--\,8) is honored exactly: the plugin is placed
  in that slot even if this leaves empty slots (gaps) in the FX chain. On
  older REAPER versions, or with that option disabled, the fixed position is
  only treated as a maximum: if the chain contains fewer FX, the plugin is
  inserted earlier (in a higher slot) instead of leaving gaps. The
  \emph{first} and \emph{last} positions are not affected by this.
\item \textbf{Edit Mapping}: Opens the mapping editor for this strip, see
  \hyperref[stripmapping]{next section}.
\item \textbf{Save}: Write the strip (plugin, abbreviation and \vpot
  mapping) to a file that you choose in a file dialog. The file name is
  pre-filled with the strip's abbreviation.
\item \textbf{Load}: Replace the strip with the first strip found in a file
  that you choose. The current contents of this row are overwritten without
  a confirmation.
\item \textbf{Clear}: Reset the whole strip (plugin, abbreviation and \vpot
  mapping) to unassigned. \textbf{Save} and \textbf{Clear} are only enabled
  for strips that have a plugin; \textbf{Load} is always available.
\end{itemize}

Below the table, the buttons \textbf{Save all 16...} and
\textbf{Load all 16...} work on the complete set: \emph{Save all} writes
every assigned strip to one file, and \emph{Load all} \textbf{replaces all
16 slots} with the strips found in the file --- slots that the file does not
contain become unassigned. A file that does not contain a single valid strip
is rejected, so an accidental load of a wrong file cannot wipe the set.

\subsubsection{Mapping Editor}\label{stripmapping}
The ``Edit Mapping'' button of a strip opens the
\hyperref[F:Screenshot_ChannelStrip_Mapping]{VPOT mapping editor} (see
Figure~\ref{F:Screenshot_ChannelStrip_Mapping}). It shows 16 rows
(VPOT 1\,--\,8 and Shift 1\,--\,Shift 8), one for each \vpot of the strip.

\pikFigure[0.6][9.0cm][-2.2mm]{Screenshot_ChannelStrip_Mapping}{VPOT mapping
  editor (\texttt{Edit Mapping})}{}{}

For each row you select the parameter of the strip's plugin that the \vpot
should control in the \textbf{Parameter} column. The parameter can be picked
directly from the combo box, or you can select the row and then change the
corresponding value in the plugin's own user interface; the parameter that
changes is assigned to the selected row. The \textbf{Name} column holds a
short display label (up to 6 characters) that is shown on the MCU; it is
filled automatically from the parameter name and can be edited afterwards.

The mapping of a strip is applied to every unit that is assigned to that
strip. If the plugin is not present on the selected track, the editor adds a
temporary instance for the purpose of learning the mapping and removes it
again when the editor is closed.

\subsubsection{Saving and Loading Strips}\label{stripsave}
Strips can be exchanged as files, either single strips (the \textbf{Save}
and \textbf{Load} buttons of the table) or the complete set of all 16
strips (the \textbf{Save all 16...} and \textbf{Load all 16...} buttons
below the table).

The file dialog always opens in the Channel Strip Maps directory of the
user location:

\begin{itemize}
\item \reaper on Windows: \texttt{<Documents>\textbackslash Reaper\textbackslash MCU\textbackslash ChannelStripMaps\textbackslash}
\item \reaper on macOS: \texttt{~/Library/Application Support/REAPER/MCU/ChannelStripMaps/}
\item \reaper on Linux: \texttt{~/.config/REAPER/MCU/ChannelStripMaps/}
\end{itemize}

The file format is plain XML and identical to the internal
\texttt{channelstrips.xml} file that stores the strips between \reaper
restarts: a root element that contains one entry per strip (plugin,
abbreviation, insert position and the 16 \vpot mappings). A single-strip
file simply contains one such entry, and a complete-set file contains up to
16. Because the format is readable, the files can also be edited by hand or
exchanged between systems.

\subsubsection{Commands}\label{channelstripcommands}
In the Channel Strip-Mode the \hyperref[commands]{Commands} act on the strip
FX of the pressed unit on the selected track. Press the \vpot while holding
the \control modifier; while \control is held, the display shows the command
label on the active units.

\begin{itemize}
\item \control + \vpot 1 (\emph{Float}): Toggle the floating window of the
  strip FX.
\item \control + \vpot 2 (\emph{Chain}): Toggle the FX chain window of the
  selected track.
\item \control + \vpot 3 (\emph{PlMode}): Switch to the
  \hyperref[plugmode]{FX-Mode} and select the strip FX there.
\item \control + \vpot 5 (\emph{Remove}): Remove the strip FX from the
  selected track.
\item \control + \vpot 7 / \vpot 8 (\emph{FXup} / \emph{FXdown}): Move the
  strip FX one slot up / down in the FX chain.
\end{itemize}

\subsubsection{Options}
The Channel Strip-Mode uses the same Main Options and 2nd Options as the
\hyperref[panmode]{Pan-Mode} (see Table~\ref{T:multitrack_options1} and
Table~\ref{T:multitrack_options2}).

%%% Local Variables:
%%% mode: latex
%%% TeX-master: "../mcu_klinke_manual"
%%% End:
